\section{Proof of the gap theorem}\label{sec:assembly}

We now assemble the two initializations through the same purity theorem.
The thermal and trial inputs used in Section~\ref{sec:initialization}
are proved in Appendices~\ref{sec:thermal120}--\ref{sec:trials}.
Those finite proofs are part of the argument.

\begin{proof}[Proof of Theorem~\ref{thm:main}]
Proposition~\ref{prop:transfer} identifies the physical partition
functions with moments of the self-adjoint spatial transfer operator.
They therefore satisfy the hypotheses used to define both purities in
Section~\ref{sec:bootstrap}. The Hamiltonians are self-adjoint and
non-scalar, as noted after \eqref{eq:main-gap}.

Proposition~\ref{prop:seed60} supplies
\[
 S(60,105/4)>7/8,\qquad T(120,105/4)>7/8.
\]
Apply Proposition~\ref{prop:bootstrap} with
\[
 (n_0,b_0,u_0,C)=(60,105/4,1/8,79/10).
\]
Its parameter inequalities are verified in
Section~\ref{sec:bootstrap}, and $Cu_0=79/80$.
It follows that the ground state is unique and
$\gamma_L>\frac4{105}\log(80/79)$ for every even $L\ge60$.

Proposition~\ref{prop:seed2304} supplies both purity defects at most
$1/100$ at the pair based at $(n_0,b_0)=(2304,784)$.
Apply the same theorem with $u_0=1/100$ and $C=5$.
Since $Cu_0=1/20$, it gives
$\gamma_L>\log(20)/784$ for every even $L\ge2304$.
Taking the limit inferior over even lengths proves the stated bound
on $\Delta_1$. Finally, Lemma~\ref{lem:representation} transfers the
spectral conclusions to the other spin-one matrix realization by
simultaneous similarity and its tensor powers.
\end{proof}

\subsection{An infinite-volume consequence}

The finite-volume estimate also controls local excitations in a
thermodynamic limit. We use the standard limiting argument recalled in
\cite[arXiv v4, End Matter, Definition~11 and Theorem~12]{Tasaki2025}.
A local operator on the spin chain $\Z$ acts nontrivially on only
finitely many sites. A state $\omega$ is a linear functional on these
operators with $\omega(I)=1$ and $\omega(A^*A)\ge0$ for every local
operator $A$. Define the local commutator
\[
 [H,A]=\sum_{j\in\Z}
 [\mathbf S_j\cdot\mathbf S_{j+1},A].
\]
Only finitely many summands are nonzero.

\begin{corollary}\label{cor:local-limit}
For the self-adjoint Hamiltonians \eqref{eq:main-H}, the normalized
ground states of the even periodic rings have a subsequence, as
$L\to\infty$, converging on every local operator to a state $\omega$.
Every such limit satisfies
\begin{equation}\label{eq:local-limit-gap}
 \omega(A^*[H,A])\ge\frac{\log20}{784}
 \bigl(\omega(A^*A)-|\omega(A)|^2\bigr)
\end{equation}
for every local operator $A$.
\end{corollary}

\begin{proof}
Place the ring of even length $L$ on
$\{-L/2,\ldots,L/2-1\}$, with its closing bond at the endpoints.
Compactness of the density matrices on each finite interval and a
diagonal subsequence give convergence on all local operators. The
compatible limits define a positive normalized functional $\omega$.
For $L\ge2304$, Theorem~\ref{thm:main} gives a unique ground state,
with normalized vector $\psi_L$, and a gap greater than
$\delta=\log(20)/784$.
Writing $\omega_L(A)=\langle\psi_L,A\psi_L\rangle$, the spectral
inequality on the orthogonal complement of $\psi_L$ gives
\[
 \omega_L(A^*[H_L,A])\ge\delta
 \bigl(\omega_L(A^*A)-|\omega_L(A)|^2\bigr).
\]
For fixed local $A$ and sufficiently large $L$, the closing bond commutes
with $A$ and $[H_L,A]=[H,A]$. Passing to the local limit proves the claim.
\end{proof}

In particular, $\omega$ is a ground state: no local perturbation lowers
its energy. For $\omega(A)=0$, inequality~\eqref{eq:local-limit-gap}
is the locally unique gapped property of
\cite[arXiv v4, Definition~11]{Tasaki2025}, with gap at least
$\log(20)/784$. The argument does not identify the limits of different
subsequences or establish uniqueness among all infinite-volume ground
states.
